\chapter{Balloon Valvuloplasty}

\section*{Overview}
Balloon valvuloplasty is a catheter-based procedure that widens a narrowed heart valve by inflating a balloon across the stenotic orifice. It is most commonly used for mitral and aortic valve stenosis. The exact approach, setting, and preparation depend on the indication, anatomy, and the patient's health.

\section*{Alternative Names}
Balloon valvotomy, percutaneous balloon valvuloplasty, mitral balloon commissurotomy, aortic balloon valvuloplasty

\section*{Principle}
A balloon catheter is advanced across the stenotic valve and inflated to high pressure, splitting fused commissures and stretching the valve leaflets. This increases the effective orifice area and reduces the transvalvular pressure gradient. The result improves forward blood flow and relieves symptoms of obstruction.

\section*{History}
Rubio and Limón Lason performed the first balloon valvuloplasty for pulmonic stenosis in 1953, and Inoue developed the modern mitral balloon technique in 1982. It became an alternative to open surgical commissurotomy in selected patients.

\section*{Why It Is Done}
Performed for symptomatic severe mitral stenosis with favorable valve morphology, and for severe aortic stenosis in patients who are poor surgical candidates, such as bridging therapy before surgical or transcatheter valve replacement, or for critical aortic stenosis in pregnancy. In aortic stenosis it provides only temporary symptomatic relief.

\section*{Is This a Primary or Secondary Test or Procedure}
Primary treatment for selected patients with mitral stenosis; a palliative or bridge procedure for aortic stenosis in those ineligible for valve replacement.

\section*{How It Is Performed}
Performed in the cardiac catheterization laboratory under conscious sedation or general anesthesia. Access is obtained via the femoral vein (mitral, using transseptal puncture) or femoral artery (aortic). Under fluoroscopic and echocardiographic guidance, the balloon is positioned across the valve, inflated, and withdrawn.

\section*{Preparation}
Transthoracic and often transesophageal echocardiography for valve scoring and left atrial thrombus exclusion. Blood typing and crossmatch, and withholding of anticoagulants as directed.

\section*{Contraindications and Precautions}
Contraindicated with severe valvular regurgitation, left atrial thrombus, or severe mitral subvalvular and leaflet calcification (for mitral). Precautions include moderate-to-severe aortic regurgitation and coexisting coronary artery disease (for aortic).

\section*{Risks}
Valvular regurgitation, systemic embolism from calcium or thrombus, vascular access complications, cardiac perforation, and arrhythmias. Restenosis is common, particularly after aortic balloon valvuloplasty.

\section*{Aftercare and Recovery}
Patients are monitored in a cardiac observation unit with echocardiography to assess valve area and regurgitation. Most are discharged within 24-48 hours and may resume normal activities shortly thereafter.

\section*{Outcome and Follow-up}
Success is defined by an increase in valve area of at least 50\% and a reduction in the mean gradient without significant regurgitation. Residual symptoms or high gradients after aortic valvuloplasty predict early recurrence.

\section*{Expected Results}
An effective valve orifice area near normal (mitral >2.0 cm$^2$; aortic >1.5 cm$^2$) with a low mean gradient and no more than mild regurgitation.

\section*{Follow-up}
Serial echocardiography to monitor valve area, gradient, and restenosis. Definitive valve replacement should be planned for patients with aortic stenosis, and repeat valvuloplasty or surgery may be needed for mitral restenosis.

\section*{When to Seek Medical Care}
Follow the procedure team's specific instructions. Seek urgent medical assessment for severe or worsening pain, uncontrolled bleeding, fever or signs of infection, trouble breathing, chest pain, fainting, new weakness or confusion, inability to urinate, or any rapidly worsening symptom. The appropriate warning signs depend on the procedure and the patient's medical condition.

\section*{References}
\begin{enumerate}
  \item MedlinePlus. Balloon valvuloplasty. U.S. National Library of Medicine; 2025.
  \item Current Medical Diagnosis and Treatment. McGraw-Hill; 2025.
\end{enumerate}

\clearpage
